\documentclass[11pt]{article}
\usepackage[a4paper,margin=25mm]{geometry}
\usepackage{amsmath,amssymb,mathtools}
\usepackage{fontspec,unicode-math}
\setmainfont{Latin Modern Roman}
\setsansfont{Latin Modern Sans}
\setmathfont{Latin Modern Math}
\usepackage{xurl,hyperref}
\hypersetup{colorlinks=true,linkcolor=blue,urlcolor=blue,pdftitle={Problem statement}}
\newcommand{\sref}[2]{\hyperlink{source:#1:#2}{[S#2]}}
\newcommand{\cataloguescope}[1]{\paragraph{Mathematical question.}#1}
\begin{document}
\section{Size of the instability locus}\label{statement}

\subsection{Definitions and mathematical statement}

\paragraph{Shared notation.}$\mathbb N=\{1,2,\ldots\}$, and $\mathbb Z,\mathbb Q,\mathbb R,\mathbb C$ denote the integers, rationals, reals and complex numbers. Any other conventions are specified locally.


Let $\mathcal A$ be an analytic parameter family of rational maps of fixed degree at least two. A parameter is $J$-stable when on a neighbourhood its Julia sets move by a holomorphic motion that conjugates the dynamics: each point varies holomorphically with the parameter and the motion is injective on each Julia set. Let $Q\subset\mathcal A$ be the complementary $J$-unstable locus. Lebesgue measure and Hausdorff dimension are taken in real local coordinates on the parameter space; $\dim\mathcal A$ in the comparison means its real dimension.
\paragraph{Statement recorded in the supplied source.}
What is the Lebesgue measure of $Q$? Is $\dim_HQ=\dim_{\mathbb R}\mathcal A$? \sref{statement}{2}
\subsection{Short English statement}
How much ordinary volume does the unstable locus occupy in an analytic rational-map family, and does it have the full Hausdorff dimension of that family?
\subsection{Sources}
\begin{description}
\item[\hypertarget{source:statement:1}{S1}] ULAM.AI and the UnsolvedMath Contributors, UnsolvedMath: A Curated Collection of Open Mathematics Problems, catalogue entry. CC BY 4.0. \url{https://huggingface.co/datasets/ulamai/UnsolvedMath}.
\item[\hypertarget{source:statement:2}{S2}] Mikhail Lyubich, ``Measure and Dimension of Julia Sets'', in Ben Bielefeld and Mikhail Lyubich (eds.), Problems in Holomorphic Dynamics (1992), Problem 6, pp. 27; arXiv:math/9205209v1. \url{https://arxiv.org/pdf/math/9205209}.
\end{description}
\label{end:statement}
\end{document}
